\textbf{Tasks and data.} Two regimes of the two-scale system: F20\_c10 (primary; strong scale separation, $c{=}10$) and F20\_c4 ($c{=}4$, weaker separation); $K{=}8$, $J{=}32$, $h{=}1$, $b{=}10$, $F{=}20$. The full model is integrated with RK4 at step 0.0025 (a logged sanity run, \texttt{experiments/precheck\_dt.py}: at $c{=}10$, the mean and the variance of $X$ over 40 trajectories and a 20-unit window each differed by less than 1\% from a run with half the step from the same initial states) and sampled every 0.005. For seed $s$ we generate independent trajectory sets after a 20-unit spin-up: 16 training chains of 25 time units (thinned by 4 for fitting), 4 validation chains (MLP early stopping only), 10 test chains of 30 units (forecast initial conditions every 3 units, offline $R^2$). Climate is scored against one fixed reference run (40 chains $\times$ 100 units, same for every system and seed). As a noise floor, an independent truth run of the same size (seed-dependent) is scored against that reference (``Truth (indep.\ run)'').

\textbf{Protocol.} Systems: no closure, polynomial (degree 4), MLP ($1$--$32$--$32$--$1$, local input), AR(1) stochastic (degree-4 polynomial plus noise); all reimplemented here, with no external code. Five seeds (0--4) for the main comparison and three (0--2) for ablations; a seed changes the training, validation and test trajectories, the MLP initialisation and batch order, and all noise streams. Rollouts: closed-model step 0.005 for every system, 10-member forecast ensembles, 20 free runs of 400 units with 10 units discarded, $X$ stored every 0.05. Hyperparameters were fixed before any logged run and not tuned (no tuning budget used); the polynomial degree and the MLP input width are examined only as ablations on the test protocol (which has no separate selection split), so the ablation rows are not a model selection. Tests: \texttt{rh compare} reports Welch and paired $t$ tests between per-seed values; we quote the Welch $p$ without multiple-comparison correction, with $n{=}5$ per system.

\textbf{Hardware and runtime.} A shared macOS CPU machine, 2 threads, two experiment processes at a time; 99 runs in the registry for the study (including one smoke-test sanity run), plus the step-size pre-check logged afterwards as a second sanity run, with a median in-process time of about 10 s for non-neural systems and a median of about 44 s for MLP runs. Two MLP runs of the main group (seed 1 at $c{=}10$ and seed 4 at $c{=}4$) took 390 s and 549 s in-process (394 s and 552 s as timed by the harness) because of contention on the shared machine, longer than the 6-minute limit we aimed for; their results are included. The registry also holds 24 superseded rows of the forcing-shift ablation: its 18 runs were first logged with shift-specific task labels, which the harness check does not accept, and the six AR(1) ones were logged a second time under an interim name; all 24 rows were superseded (\texttt{rh supersede}) and the 18 runs were executed again with the same commands and seeds under the final names. All metrics of the superseded rows other than the runtime are identical to those of the final rows. Code, per-run logs and the run registry (\texttt{results/runs.jsonl}, which holds the metrics of every run) are in the repository; the per-run JSON files and the generated data are not committed. \texttt{experiments/run\_all.sh} lists the commands of the pre-check and of all main, floor and ablation runs under their final names.
